\documentclass[11pt]{article}

\usepackage[margin=1in]{geometry}
\usepackage[T1]{fontenc}
\usepackage[english]{babel}
\usepackage{amsmath,amsthm,amssymb}
\usepackage{array}
\usepackage{mathtools}
\usepackage{stmaryrd}
\usepackage{nicefrac}
\usepackage{xcolor}
\usepackage{etoolbox}
\usepackage{infer}

\providecommand{\Form}{\Phi}
\usepackage{paper}
\providecommand{\hoellora}{\textsf{H$\mbox{}_2$O$\ell$}\xspace}
\providecommand{\HOL}{\textsf{HOL}\xspace}
\providecommand{\todo}[2][red]{}
\providecommand{\gb}[1]{}\providecommand{\ma}[1]{}\providecommand{\dd}[1]{}
\renewcommand*{\InferTagFormat}[1]{\textsc{\small{[#1]}}}
\hypersetup{colorlinks=true,linkcolor=blue,urlcolor=blue,citecolor=blue}

% notation
\newcommand{\iHL}{\ensuremath{\mathsf{IL}}\xspace}   % name of the object logic
\newcommand{\iat}{\mathbin{\ast}}                 % independence atom  E * F
\newcommand{\W}{W}
\newcommand{\wn}[1]{\tfrac{w_{#1}}{\W}}
\newcommand{\Upd}[2]{U_{#2}(#1)}
\newcommand{\Updf}{U^{\mathsf{fed}}}
\newcommand{\smpl}{\SAMPLE}
\newcommand{\Bsol}[1]{B_{#1}}
\newcommand{\grad}{\nabla\! f}
\newcommand{\stp}{\mathrel{;}}
\newcommand{\munit}{\mathsf{unit}}
\newcommand{\Law}{\mathcal L}
\newcommand{\old}[1]{\mathsf{old}(#1)}
\providecommand{\dsg}[1]{\llbracket #1 \rrbracket}
\newcommand{\evat}[2]{\Sem{#1}_{#2}}
\newcommand{\restr}[2]{#1|_{#2}}                  % restriction / conditioning  A|_b
\newcommand{\mass}[1]{\lvert #1\rvert}
\newcommand{\solo}{\mathsf{solo}}
\newcommand{\fed}{\mathsf{fed}}

\theoremstyle{plain}\newtheorem{lemma}{Lemma}\newtheorem{theorem}{Theorem}\newtheorem{corollary}{Corollary}
\theoremstyle{definition}\newtheorem{definition}{Definition}

\begin{document}

\section{An independence logic and its embedding}
\label{sec:il}

We use \hoellora\ to give semantics
and a soundness proof to a new logic, which we will refer to as \iHL, tailored to 
the propagation of \emph{probabilistic independence} through an imperative program. 

We introduce an
assertion language $\{\,S \mid \form\,\}$ that pairs an independence structure
$S$ with an ordinary first-order support predicate $\form$, equip it with a
Hoare-style proof system over a family of imperative programs, and embed it into
\hoellora. Soundness of the object logic 
then follows by composing the embedding
with the soundness of \hoellora. We also show that the derived
system helps us prove the variance-reduction guarantee of fedSGD.


\subsection{The logic \iHL}

\paragraph{Syntax.}
The logic will let us prove judgements over a family of programs written by the
simple imperative programming language generated by this grammar:
\[
  c \ ::=\ x\leftarrow E \ \mid\ x\smpl d \ \mid\ c\stp c
        \ \mid\ \kwd{for}\ k{=}a\ \kwd{to}\ b\ \kwd{do}\ c
        \ \mid\ \kwd{if}\ b\ \kwd{then}\ c\ \kwd{else}\ c\,.
\]
As we mentioned, assertions in this logic will take the form $\{\,S \mid \form\,\}$.
Support formulas $\form$ are just first-order formulas over expressions. Whereas, 
independence structures $S$ are those generated by the following grammar:
\[
  S \ ::=\ \top \ \mid\ E \iat F \ \mid\ \restr{S}{b} \ \mid\ S\land S\,.
\]
where $E,F$ range over expressions over program variables, and $b$ over boolean
expressions. The atom $E\iat F$ asserts independence of the two expressions; $\restr{S}{b}$ is
\emph{$S$ conditioned on $b$}, i.e.\ independence \emph{conditional} on the guard
$b$ (it holds on the fibre where $b$ is true); and $\land$ conjoins structures.

\paragraph{Semantics.}
\iHL\ is interpreted directly in \hoellora: commands as terms of its programming
language, and the satisfaction relations below as formulas of its ambient \HOL. 
Commands are interpreted as follows:
\[
  \begin{array}{@{}l@{}}
  \Sem{x\leftarrow E}=\lam{m}{\ret{m[x\mapsto\evat{E}{m}]}},\\[0.3ex]
  \Sem{x\smpl d}=\lam{m}{\mletin{v}{\Sem{d}\,m}{\ret{m[x\mapsto v]}}},\\[0.3ex]
  \Sem{c\stp c'}=\lam{m}{\mletin{m'}{\Sem{c}\,m}{\Sem{c'}\,m'}},\\[0.3ex]
  \Sem{\kwd{if}\ b\ \kwd{then}\ c\ \kwd{else}\ d}=\lam{m}{\cif{\evat{b}{m}}{\Sem{c}\,m}{\Sem{d}\,m}},\\[0.3ex]
  \Sem{\kwd{for}\ k{=}a\ \kwd{to}\ b\ \kwd{do}\ c}=\lam{m}{\op[iter]\,(b{-}a{+}1)\,(\lam{k,\mu}{\mletin{m'}{\mu}{\Sem{c}\,m'}})\,(\ret m)}.
  \end{array}
\]

An independence structure is interpreted against a coupling $\dist$. The relation
$\dist\models S$, read ``the coupling $\dist$ has independence structure $S$'', is
defined by induction on $S$:
\begin{align*}
  \dist \models \top\ &\text{always;}\\
  \dist \models E\iat F\ &\Longleftrightarrow\
    \mass{\dist}\cdot\mbind{\dist}{\lam{\hvec m}{h(\evat{E}{\hvec m},\evat{F}{\hvec m})}}\\
  &\hphantom{\Longleftrightarrow\ }\quad
    {}=\mbind{\dist}{\lam{\hvec m_1}{\mbind{\dist}{\lam{\hvec m_2}{h(\evat{E}{\hvec m_1},\evat{F}{\hvec m_2})}}}}
    \quad\text{for all }h;\\
  \dist \models S_1\land S_2\ &\Longleftrightarrow\ \dist\models S_1\ \text{and}\ \dist\models S_2;\\
  \dist \models \restr{S}{b}\ &\Longleftrightarrow\ \restr{\dist}{b}\models S.
\end{align*}
Here $\mass{\dist}$ is the total mass of $\dist$ and, for a boolean guard $b$,
$\restr{\dist}{b}$ is the subdistribution of $\dist$ supported on
$\{\hvec m\mid\evat{b}{\hvec m}\}$. The clause for $E\iat F$ is the independence of
$E$ and $F$ under $\dist$ in scale-invariant form: the factor $\mass{\dist}$ keeps it
well defined on the subdistributions $\restr{\dist}{b}$ produced by conditioning (for
couplings of mass $1$ it drops out, leaving the usual factorisation). The clause for
$\restr{S}{b}$ thus reads independence \emph{conditional} on $b$; when
$\Prob{\dist}{b}=0$ it holds vacuously, as $\restr{\dist}{b}=0$.

An assertion is satisfied by a coupling $\dist$, written
$\dist\models\{\,S\mid\form\,\}$, if $\dist\models S$ and the support formula holds
at every point of the support, $\fa{\hvec m\in\supp(\dist)}{\hvec m\models\form}$.

\begin{definition}[Semantic validity, $\models$]\label{def:il-valid}
A family of marginals $(\dist_i)_{i\in I}$ satisfies an assertion, written
$(\dist_i)_{i\in I}\models\{\,S\mid\form\,\}$, if some coupling of the family does:
\[
  (\dist_i)_{i\in I}\models\{\,S\mid\form\,\}\ \defsym\
  \exists\dist.\ \dist\ \text{couples}\ (\dist_i)_i \land \dist\models\{\,S\mid\form\,\}.
\]
Then
\[
  \models \{\,S\mid\form\,\}\ [\,i\colon c_i\,]_{i\in I}\ \{\,S'\mid\form'\,\}
  \ \defsym\
  \fa{(\dist_i)_{i\in I}}{\ (\dist_i)_i\models\{\,S\mid\form\,\} \Rightarrow (\mletin{m}{\dist_i}{\Sem{c_i}\,m})_i\models\{\,S'\mid\form'\,\}\ }.
\]
\end{definition}

\paragraph{Proof system.}
{\small
\begin{gather*}
  \Infer[Assign]
  { \judge \{\, S[x/\hvec E] \mid \form[x/\hvec E] \,\}\ [\,i\colon x\leftarrow E_i\,]_{i\in I}\ \{\, S\mid\form \,\} }
  { }
  \\[1.6ex]
  \Infer[Sample]
  { \judge\ \{\, S \mid \fa{\hvec v\in\supp(\dist)}{\form[x/\hvec v]} \,\}\ [\,i\colon x\smpl d_i\,]_{i\in I}\ \{\, S'\mid\form \,\} }
  { \begin{array}[c]{@{}c@{}}
      \dist \mathrel{\triangleright} d_1 \mathbin{\&} \cdots \mathbin{\&} d_n\\[0.4ex]
      \fa{\disttwo\models S}{\ \mbind{\disttwo}{\lam{\hvec m}{\mbind{(\dist\,\hvec m)}{\lam{\hvec v}{\munit(\hvec m[x/\hvec v])}}}}\ \models\ S'\ }
    \end{array} }
  \\[1.6ex]
  \Infer[Sample$_2$]
  { \judge\ \{ \textstyle\bigwedge_{i\ne j\in J} A_i \iat A_j \mid \fa{\hvec v\in\supp(\dist)}{\form[x/\hvec v]} \}\ [\,i\in I\colon x\smpl d_i\,]\ \{ \textstyle\bigwedge_{i\ne j\in J} (A_i, x_i) \iat (A_j, x_j) \mid \form \} }
  { \begin{array}[c]{@{}c@{}}
      J\subseteq I \quad \dist \mathrel{\triangleright} d_1 \mathbin{\&} \cdots \mathbin{\&} d_n\\[0.4ex]
      \fa{i\ne j\in J}{\ \exists e_i,e_j.\ \fa{\hvec m}{\ \pi_{\{i,j\}}(\dist\,\hvec m)=e_i\otimes e_j\ }}
    \end{array} }
  \\[1.6ex]
  \Infer[Conseq]
  { \judge \{\,S'\mid\form'\,\}\ [\,i\colon c_i\,]_{i\in I}\ \{\,S''\mid\form''\,\} }
  { S'\vdash S & \form'\Rightarrow\form
    & \judge \{\,S\mid\form\,\}\ [\,i\colon c_i\,]_{i\in I}\ \{\,T\mid\psi\,\}
    & T\vdash S'' & \psi\Rightarrow\form'' }
  \\[1.6ex]
  \Infer[Seq]
  { \judge \{\,S\mid\form\,\}\ [\,i\colon c_i\stp c_i'\,]_{i\in I}\ \{\,S''\mid\form''\,\} }
  { \judge \{\,S\mid\form\,\}\ [\,i\colon c_i\,]_{i\in I}\ \{\,S'\mid\form'\,\}
    & \judge \{\,S'\mid\form'\,\}\ [\,i\colon c_i'\,]_{i\in I}\ \{\,S''\mid\form''\,\} }
  \\[1.6ex]
  \Infer[For]
  { \judge \{\,\mathcal J(a)\,\}\ [\,i\colon \kwd{for}\ k{=}a\ \kwd{to}\ b\ \kwd{do}\ c_i\,]_{i\in I}\ \{\,\mathcal J(b{+}1)\,\} }
  { \fa{k\in\{a,\dots,b\}}{\ \judge \{\,\mathcal J(k)\,\}\ [\,i\colon c_i\,]_{i\in I}\ \{\,\mathcal J(k{+}1)\,\}\ } }
  \\[1.6ex]
  \Infer[If]
  { \judge \{\restr{S}{b}\land\restr{S}{\lnot b}\mid\form\}\ [\,i\colon\kwd{if}\ b_i\ \kwd{then}\ c_i\ \kwd{else}\ d_i\,]_{i\in I}\ \{ \restr{T}{p}\land\restr{U}{\lnot p} \mid (p\Rightarrow\form^c)\land(\lnot p\Rightarrow\form^d) \} }
  { \begin{array}[c]{@{}c@{}}
      \form\vdash b_1{=}\cdots{=}b_n{=:}b\\[0.4ex]
      \judge \{\,S\mid\form\land b\,\}\,[\,i\colon c_i\,]_{i\in I}\,\{\,T\mid\form^c\,\}
      \quad \judgeHOL\ \form^c\Rightarrow p\\[0.4ex]
      \judge \{\,S\mid\form\land\lnot b\,\}\,[\,i\colon d_i\,]_{i\in I}\,\{\,U\mid\form^d\,\}
      \quad \judgeHOL\ \form^d\Rightarrow\lnot p
    \end{array} }
\end{gather*}
}

In rules \textsc{[Sample]} and \textsc{[Sample$_2$]}, $\dist$ is a store-indexed
coupling of the samplers, i.e.\ a function from tuples $\hvec m=(m_i)_{i\in I}$ of
the stores on which the $\Sem{c_i}$ act to distributions over tuples of sampled
values, and
\[
  \dist \mathrel{\triangleright} d_1 \mathbin{\&} \cdots \mathbin{\&} d_n
  \ \defsym\
  \fa{\hvec m}{\ \fa{i\in I}{\ \pi_i(\dist\,\hvec m)=\evat{d_i}{m_i}\ }}
\]
with $\pi_i$ the $i$-th marginal. In support formulas, $\dist$ stands for
$\dist\,\hvec m$ at the current store tuple $\hvec m$. In rule
\textsc{[Sample$_2$]},
$e\otimes e'\defsym\mbind{e}{\lam{u}{\mbind{e'}{\lam{w}{\munit(u,w)}}}}$ denotes
the product of two distributions.


\paragraph{Assertion logic.}
We present some structural facts about $S$ which may be useful to use with the 
program logic through \textsc{[Conseq]}.
{\small
\begin{align*}
  \textsc{[Affine]}\ \ & S \land (E\iat F) \ \vdash\ S, &
  \textsc{[Map]}\ \ & \hvec E \iat \hvec F \ \vdash\ f(\hvec E)\iat g(\hvec F),\\
  \textsc{[Cond-$\land$]}\ \ & \restr{(S_1\land S_2)}{b} \ \mathrel{\dashv\vdash}\ \restr{S_1}{b}\land\restr{S_2}{b}, &
  \textsc{[Cond-Cond]}\ \ & \restr{\bigl(\restr{S}{b'}\bigr)}{b} \ \mathrel{\dashv\vdash}\ \restr{S}{\,b\land b'},\\
  \textsc{[Const]}\ \ & \top \ \vdash\ E \iat c \quad (c\ \text{constant}).
\end{align*}
}


\subsection{Embedding into \hoellora}

\label{sec:il-embed}

Writing $\dsg{\{S\mid\form\}}(\hvec\dist)$ for the \HOL\ formula
$(\dist_i)_i\models\{\,S\mid\form\,\}$ of \S\ref{sec:il}, we interpret an \iHL\
triple as follows:
\[
  \dsg{\ \judge \{\,S\mid\form\,\}\ [\,i\colon c_i\,]_{i\in I}\ \{\,S'\mid\form'\,\}\ }
  \ \defsym\
  (\dist_i\ofty\tydist\,\sigma)_{i}\ \bigm|\ \dsg{\{S\mid\form\}}(\dist_i)_i
  \ \judge\ \HMAP{\Ix{i}\, \mletin{m}{\dist_i}{\Sem{c_i}\,m}}_{i\in I}
  \ \ASRT{\dsg{\{S'\mid\form'\}}}.
\]

\begin{theorem}[Embedding of \iHL\ in \hoellora]\label{thm:il-emb}
If $\judge\{\,S\mid\form\,\}\ [\,i\colon c_i\,]_{i\in I}\ \{\,S'\mid\form'\,\}$,
then $\dsg{\,\judge\{\,S\mid\form\,\}\ [\,i\colon c_i\,]_{i\in I}\ \{\,S'\mid\form'\,\}\,}$
is derivable in \hoellora.
\end{theorem}

The proof is presented in Appendix~\ref{app:il}.

\begin{corollary}[Soundness of \iHL]\label{cor:il-sound}
If $\judge\{\,S\mid\form\,\}\ [\,i\colon c_i\,]_{i\in I}\ \{\,S'\mid\form'\,\}$
then $\models\{\,S\mid\form\,\}\ [\,i\colon c_i\,]_{i\in I}\ \{\,S'\mid\form'\,\}$.
\end{corollary}


\subsection{Application: variance reduction in federated SGD}

\label{sec:il-fedsgd}

Federated SGD has $n$ parties collaboratively train a shared model: at each
round, every party computes a gradient on a minibatch drawn from its own local
dataset $D_i$, and a central server aggregates these per-party gradients with
fixed weights $w_1,\dots,w_n$. When the parties' datasets are independent, the
aggregate enjoys a variance reduction over any single solo run of SGD. Let
$\W=\sum_j w_j$, let the affine gradient be $\grad(\theta,x)=\theta-x$ (square
loss), and write the batch update $\Upd{\theta}{B}\defsym\theta-\eta\,\tfrac1m\sum_{x\in B}\grad(\theta,x)$.
The index set is $I=\{1,\dots,n,\,n{+}1\}$: for $i\le n$ the $i$-th program runs
$\solo$ on the local dataset $D_i$; the $(n{+}1)$-th program runs $\fed$ on
$\hvec D=(D_1,\dots,D_n)$. There are $R$ rounds.
\begin{center}
\small
\begin{tabular}{@{}p{0.46\textwidth}p{0.5\textwidth}@{}}
\begin{minipage}[t]{\linewidth}
\begin{tabbing}
\ \=\ \=\kill
$\solo(\theta_{\mathsf{init}},D)$:\\
\>$\theta\leftarrow\theta_{\mathsf{init}}$;\\
\>$\kwd{for}\ k{=}1\ \kwd{to}\ R\ \kwd{do}$\\
\>\>$B \smpl \op[minibatch]\ m\ D$;\\
\>\>$\theta \leftarrow \theta-\eta\,\tfrac1m\textstyle\sum_{x\in B}\grad(\theta,x)$
\end{tabbing}
\end{minipage}
&
\begin{minipage}[t]{\linewidth}
\begin{tabbing}
\ \=\ \=\kill
$\fed(\theta_{\mathsf{init}},\hvec D,\hvec w)$:\\
\>$\theta\leftarrow\theta_{\mathsf{init}}$;\\
\>$\kwd{for}\ k{=}1\ \kwd{to}\ R\ \kwd{do}$\\
\>\>$B \smpl \op[minibatchN]\ m\ \hvec D$\quad(\,$B=(B_1,\dots,B_n)$\,);\\
\>\>$\theta \leftarrow \theta-\eta\,\textstyle\sum_{k\le n}\wn{k}\tfrac1m\sum_{x\in B_k}\grad(\theta,x)$
\end{tabbing}
\end{minipage}
\end{tabular}
\end{center}
We write $\theta_i$ ($i\le n$), $\theta_{n+1}$ for the models, $\Bsol{i}$ for the
batch of program $i\le n$, and $(\Bsol{1}^{\mathsf f},\dots,\Bsol{n}^{\mathsf f})$
for the batch tuple of program $n{+}1$. The assertion we want to reach is
\[
  \mathcal I \ \defsym\ \Bigl\{\ \underbrace{\textstyle\bigwedge_{i\ne j\le n}\theta_i\iat\theta_j}_{S}\ \Bigm|\ \underbrace{\textstyle\sum_{i\le n}\wn{i}\,\theta_i=\theta_{n+1}}_{\form}\ \Bigr\}\,:
\]
the local models are pairwise independent, and the federated model is their
weighted mean. The body of each program is $c_i^{\mathsf{smp}}\stp c_i^{\mathsf{upd}}$, where the
sampling command is $c_i^{\mathsf{smp}}=(B\smpl d_i)$ with
$d_i=\op[minibatch]\,m\,D_i$ for $i\le n$ and $d_{n+1}=\op[minibatchN]\,m\,\hvec D$,
and the update command is $c_i^{\mathsf{upd}}=(\theta\leftarrow\Upd{\theta}{B})$,
that is $\theta\leftarrow\Upd{\theta_i}{\Bsol{i}}$ for $i\le n$ and
$\theta\leftarrow\Updf(\theta_{n+1})$ for the federated program, with
$\Updf(\theta_{n+1})\defsym\theta_{n+1}-\eta\sum_{k\le n}\wn{k}\tfrac1m\sum_{x\in \Bsol{k}^{\mathsf f}}\grad(\theta_{n+1},x)$.

We show $\judge\{\,\top\mid\top\,\}\ \hvec c\ \mathcal I$, where
$\hvec c=[\,i\colon c_i\,]_{i\in I}$ is the family with
$c_i=\solo(\theta_{\mathsf{init}},D_i)$ for $i\le n$ and
$c_{n+1}=\fed(\theta_{\mathsf{init}},\hvec D,\hvec w)$. Applying \textsc{[Seq]} to
split the initialisation from the loop, the initialisation establishes
\[
  \judge\ \{\,\top\mid\top\,\}\ [\,i\colon\theta\leftarrow\theta_{\mathsf{init}}\,]_{i\in I}\ \mathcal I,
\]
by \textsc{[Assign]} with \textsc{[Conseq]}: the substituted precondition
$\mathcal I[\theta/\theta_{\mathsf{init}}]$ has structure
$\bigwedge_{i\ne j\le n}\theta_{\mathsf{init}}\iat\theta_{\mathsf{init}}$, entailed
by $\top$ through \textsc{[Const]}, and support
$\sum_{i\le n}\wn{i}\,\theta_{\mathsf{init}}=\theta_{\mathsf{init}}$, valid since
$\sum_i w_i=\W$. Then \textsc{[For]} with the round-independent invariant
$\mathcal J(k)\equiv\mathcal I$ reduces the goal to the single-iteration
obligation
\begin{equation*}
  \judge\ \mathcal I\ \ [\,i\colon c_i^{\mathsf{smp}}\stp c_i^{\mathsf{upd}}\,]_{i\in I}\ \ \mathcal I. \tag{$\ast$}
\end{equation*}

For the sampling command,
\begin{equation*}
  \judge\ \mathcal I\ \ [\,i\colon c_i^{\mathsf{smp}}\,]_{i\in I}\ \
  \bigl\{\, S_1 \bigm| \form\land\textstyle\bigwedge_{k\le n}\Bsol{k}^{\mathsf f}=\Bsol{k}\,\bigr\}, \tag{1}
\end{equation*}
with $S_1=\bigwedge_{i\ne j\le n}(\theta_i,\Bsol{i})\iat\{\theta_j,\Bsol{j}\}$, by
\textsc{[Sample$_2$]} with $J=\{1,\dots,n\}$ and 
$A_i=\theta_i$, instantiated at the coupling $\dist$ that draws
$\Bsol{1},\dots,\Bsol{n}$ from $\bigotimes_{k\le n}\op[minibatch](D_k)$ and
returns $(\Bsol{1},\dots,\Bsol{n})$ as program $n{+}1$'s batch tuple. This
$\dist$ meets the premises of \textsc{[Sample$_2$]} and on its entire support satisfies
$\Bsol{k}^{\mathsf f}=\Bsol{k}$ for $k\le n$. Appending each fresh $\Bsol{i}$ to
its group $A_i$ produces the structure $S_1$; and, taking postcondition support
formula $\form\land\bigwedge_{k\le n}\Bsol{k}^{\mathsf f}=\Bsol{k}$, the
precondition required by \textsc{[Sample$_2$]} is
$\fa{\hvec v\in\supp\dist}{\bigl(\form\land\textstyle\bigwedge_k v^{\mathsf f}_k=v_k\bigr)}$,
where $\hvec v$ ranges over the jointly sampled batch tuples. Distributing the
quantifier over the conjunction, this is
$\bigl(\fa{\hvec v\in\supp\dist}{\form}\bigr)\land\bigl(\fa{\hvec v\in\supp\dist}{\textstyle\bigwedge_k v^{\mathsf f}_k=v_k}\bigr)$:
the first conjunct is $\form$, since $\form$ contains no batch variable and so does
not depend on $\hvec v$; the second is valid, since every $\hvec v\in\supp\dist$
satisfies $v^{\mathsf f}_k=v_k$. The precondition is thus $\form$, so we get the
invariant $\mathcal I$.

For the update command,
\begin{equation*}
  \judge\ \bigl\{\, S_2 \bigm| \psi \,\bigr\}\ \ [\,i\colon c_i^{\mathsf{upd}}\,]_{i\in I}\ \ \mathcal I, \tag{2}
\end{equation*}
where $S_2=\bigwedge_{i\ne j\le n}\Upd{\theta_i}{\Bsol{i}}\iat\Upd{\theta_j}{\Bsol{j}}$
and $\psi=\bigl(\sum_{i\le n}\wn{i}\,\Upd{\theta_i}{\Bsol{i}}=\Updf(\theta_{n+1})\bigr)$,
directly by \textsc{[Assign]}, since $\{\,S_2\mid\psi\,\}$ is $\mathcal I$ with each
$\theta_i$ replaced by the update it receives.

\begin{lemma}[Linearity of the aggregate]\label{lem:il-c2}
If $\sum_{i\le n}\wn{i}\theta_i=\theta_{n+1}$ and $\Bsol{k}^{\mathsf f}=\Bsol{k}$
for all $k\le n$, then $\sum_{i\le n}\wn{i}\,\Upd{\theta_i}{\Bsol{i}}=\Updf(\theta_{n+1})$.
\end{lemma}
\begin{proof}
With $\grad(\theta,x)=\theta-x$ we have $\Upd{\theta}{B}=(1-\eta)\theta+\eta\,\overline B$,
$\overline B=\tfrac1m\sum_{x\in B}x$. Then
$\sum_i\wn{i}\Upd{\theta_i}{\Bsol{i}}=(1-\eta)\theta_{n+1}+\eta\sum_i\wn{i}\overline{\Bsol{i}}$
using $\sum_i\wn{i}\theta_i=\theta_{n+1}$, while
$\Updf(\theta_{n+1})=(1-\eta)\theta_{n+1}+\eta\sum_k\wn{k}\overline{\Bsol{k}^{\mathsf f}}$
(as $\sum_k\wn{k}=1$); the two agree since $\Bsol{k}^{\mathsf f}=\Bsol{k}$.
\end{proof}

Composing (1) and (2) closes the argument. The postcondition of (1) entails the
precondition of (2): its structure $S_1$ gives $S_2$ by \textsc{[Map]}, because
$\Upd{\theta_i}{\Bsol{i}}$ is a function of the pair $(\theta_i,\Bsol{i})$; and
its support formula $\form\land\bigwedge_{k\le n}\Bsol{k}^{\mathsf f}=\Bsol{k}$
gives $\psi$ by Lemma~\ref{lem:il-c2}. Hence \textsc{[Conseq]} turns (1) into
$\judge\mathcal I\ [\,i\colon c_i^{\mathsf{smp}}\,]_{i\in I}\ \{\,S_2\mid\psi\,\}$,
and \textsc{[Seq]} with (2) yields the single-iteration obligation~($\ast$),
establishing $\judge\mathcal I\ \hvec c\ \mathcal I$.

With pairwise independence $S$ and the weighted-mean relation $\form$ in force at
the end, the variance reduction
$\VV(\theta_{n+1})=\sum_{i\le n}\wn{i}^{\,2}\,\VV(\theta_i)$ follows as an \HOL\
consequence.


\appendix
\section{Soundness of the embedding}

\label{app:il}

We prove Theorem~\ref{thm:il-emb} by induction on the derivation of the \iHL\
triple, showing that each rule translates to a \hoellora\ derivation of the
interpreted judgment. A substitution lemma supports the assignment case.

\begin{lemma}[Substitution]\label{lem:il-subst}
For every assertion $A$,
\[
  (\dist_i)_i\models A[x/\hvec E] \Leftrightarrow \bigl(\mletin{m}{\dist_i}{\ret{m[x\mapsto\evat{E_i}{m}]}}\bigr)_i\models A.
\]
\end{lemma}
\begin{proof}
By induction on $A$: expression substitution commutes with the pushforward along
the deterministic assignment $\lam{m}{\ret{m[x\mapsto\evat{E_i}{m}]}}$; the atoms
$E\iat F$, the conditioning $\restr{S}{b}$ (substituting into the guard as well),
and the support formula are each stable under this pushforward.
\end{proof}

\begin{proof}[Proof of Theorem~\ref{thm:il-emb}]
By induction on the derivation. The interpreted judgment has the input marginals
$(\dist_i)$ as free variables, the interpreted precondition
$\dsg{\{S\mid\form\}}(\dist_i)_i$ as its logical context $\Phi$, and the output
family $(\mletin{m}{\dist_i}{\Sem{c_i}\,m})$ as the output distributions; under $\Phi$ we derive the
postcondition $\dsg{\{S'\mid\form'\}}$ of that family, discharging each case by the
corresponding \hoellora\ rule applied to the monadic terms $\mletin{m}{\dist_i}{\Sem{c_i}\,m}$.
\begin{description}
\item[\textsc{[Assign]}]
  $\Sem{c_i}=\lam{m}{\ret{m[x\mapsto\evat{E_i}{m}]}}$ is a deterministic $\ret$,
  so by \textsc{MUnit} the output is the pushforward $\mletin{m}{\dist_i}{\Sem{c_i}\,m}$.
  By Lemma~\ref{lem:il-subst},
  $(\dist_i)_i\models\{S[x/\hvec E]\mid\form[x/\hvec E]\}\Leftrightarrow(\mletin{m}{\dist_i}{\Sem{c_i}\,m})_i\models\{S\mid\form\}$,
  which is the required implication, discharged by \textsc{Conseq}.
\item[\textsc{[Sample]}]
  $\Sem{c_i}=\lam{m}{\mletin{v}{\Sem{d_i}\,m}{\ret{m[x\mapsto v]}}}$ is a
  monadic bind, discharged by \textsc{MBind} with the coupling $\dist$ of the rule
  instance as witness, whose marginals are the samplers by the first premise.
  \textsc{MBind} establishes the structure $S'$ from its premise on the composite
  coupling, which is the second premise of \textsc{[Sample]},
  \[
    \fa{\disttwo\models S}{\ \mbind{\disttwo}{\lam{\hvec m}{\mbind{(\dist\,\hvec m)}{\lam{\hvec v}{\munit(\hvec m[x/\hvec v])}}}}\ \models\ S'};
  \]
  the support $\form$ follows from the precondition
  $\fa{\hvec v\in\supp\dist}{\form[x/\hvec v]}$, which holds at every sampled value
  and hence at every point $\hvec m[x/\hvec v]$ of the output.
\item[\textsc{[Sample$_2$]}]
  This specialises \textsc{[Sample]} to $S=\bigwedge_{i\ne j\in J}A_i\iat A_j$ and
  $S'=\bigwedge_{i\ne j\in J}(A_i,x_i)\iat(A_j,x_j)$: the sampling is the same monadic
  bind, discharged by \textsc{MBind}, so it remains to check \textsc{[Sample]}'s
  second premise: $\rho\models S'$ for every $\disttwo\models S$, where
  $\rho=\mbind{\disttwo}{\lam{\hvec m}{\mbind{(\dist\,\hvec m)}{\lam{\hvec v}{\munit(\hvec m[x/\hvec v])}}}}$.
  Fix $i\ne j\in J$. Since $\pi_{\{i,j\}}(\dist\,\hvec m)=e_i\otimes e_j$ for every
  $\hvec m$, the fresh $x_i,x_j$ are drawn from a fixed product, so under $\rho$ they
  are independent of the pre-state pair $(A_i,A_j)$ and of each other; with $A_i\iat A_j$
  from $S$, the bundles $(A_i,x_i)$ and $(A_j,x_j)$ are independent, i.e.\
  $(A_i,x_i)\iat(A_j,x_j)$. The relations of the programs in $I\setminus J$
  travel in $\form$ through $\fa{\hvec v\in\supp\dist}{}$, as in
  \textsc{[Sample]}.
\item[\textsc{[Seq]}]
  $\Sem{c_i\stp c_i'}=\lam{m}{\mletin{m'}{\Sem{c_i}\,m}{\Sem{c_i'}\,m'}}$ is Kleisli
  composition; the case follows from the two induction hypotheses and
  \textsc{MBind}.
\item[\textsc{[For]}]
  $\Sem{\kwd{for}\dots}=\lam{m}{\op[iter]\,(b{-}a{+}1)\,(\lam{k,\mu}{\mletin{m'}{\mu}{\Sem{c_i}\,m'}})\,(\ret m)}$; by
  \textsc{Iter} of \hoellora\ with invariant
  $\dsg{\mathcal J(k)}$, whose step premise is the induction hypothesis
  $\{\mathcal J(k)\}\,[\,i\colon c_i\,]\,\{\mathcal J(k{+}1)\}$.
\item[\textsc{[Conseq]}]
  \hoellora's \textsc{Conseq}. The object side-conditions translate to \HOL\
  entailments under $\dsg{\cdot}$: the structure entailments $S'\vdash S$ and
  $T\vdash S''$ give $\dsg{S'}\Rightarrow\dsg{S}$ and $\dsg{T}\Rightarrow\dsg{S''}$
  by soundness of the assertion logic (each rule \textsc{[Affine]},
  \textsc{[Map]}, \textsc{[Cond-$\land$]}, \textsc{[Cond-Cond]}, \textsc{[Const]}
  is \HOL-valid),
  while $\form'\Rightarrow\form$ and
  $\psi\Rightarrow\form''$ are \HOL\ implications on the support. \hoellora's
  \textsc{Conseq} then strengthens the precondition and weakens the postcondition.
\item[\textsc{[If]}]
  Under synchronised guards ($\form\vdash b_1{=}\cdots{=}b_n{=:}b$) the family
  denotes the branch mixture $\rho=\rho_c\uplus\rho_d$, with
  $\rho_c=\Sem{[\,i\colon c_i\,]}(\restr{\dist}{b})$ and
  $\rho_d=\Sem{[\,i\colon d_i\,]}(\restr{\dist}{\lnot b})$. Since
  $\kwd{if}\ b\ \kwd{then}\ c\ \kwd{else}\ d
  =\kwd{case}\ b\ \kwd{of}\ \mathsf{tt}\to c\mid\mathsf{ff}\to d$ (on
  $\mathsf{Bool}=1+1$), the family is discharged by \hoellora's lock-step rule
  \textsc{Case}, with guard relation $\psi(\hvec v)\defsym(v_1{=}\cdots{=}v_n)$,
  which holds by the synchronisation premise and makes every non-constant branch
  vector contradictory; the premises of the two constant vectors are the induction
  hypotheses. It remains to check the postcondition
  $\varphi\defsym\dsg{\{\restr{T}{p}\land\restr{U}{\lnot p}\mid(p\Rightarrow\form^c)\land(\lnot p\Rightarrow\form^d)\}}$
  of the output $\rho$. By the induction hypothesis
  $\rho_c\models\{\,T\mid\form^c\,\}$, so $\form^c$ holds on $\supp\rho_c$; with
  the side condition $\judgeHOL\form^c\Rightarrow p$ this gives
  $\supp\rho_c\subseteq\{\hvec m\mid\evat{p}{\hvec m}\}$, and symmetrically
  $\supp\rho_d\subseteq\{\hvec m\mid\lnot\evat{p}{\hvec m}\}$. Restriction is pointwise, so it distributes
  over the mixture and each summand is recovered exactly:
  \[
    \restr{\rho}{p}=\restr{\rho_c}{p}\uplus\restr{\rho_d}{p}=\rho_c\uplus 0=\rho_c,
    \qquad
    \restr{\rho}{\lnot p}=\rho_d.
  \]
  Hence $\rho\models\restr{T}{p}$ and $\rho\models\restr{U}{\lnot p}$ are the two
  induction hypotheses, and the support conjunct
  $(p\Rightarrow\form^c)\land(\lnot p\Rightarrow\form^d)$ holds at every point of
  $\supp\rho=\supp\rho_c\cup\supp\rho_d$, each implication being witnessed on one
  side and vacuous on the other. \textsc{Case} concludes $\varphi$ of the output.
\end{description}
This exhausts the rules, proving the theorem. Corollary~\ref{cor:il-sound} then
follows by Theorem~4.2: for the interpreted judgment its conclusion
$\Gamma\models\bigwedge\Phi\Rightarrow\phi\,\hvec s$ unfolds to
\[
  \fa{(\dist_i)_i}{\ (\dist_i)_i\models\{\,S\mid\form\,\}\Rightarrow(\mletin{m}{\dist_i}{\Sem{c_i}\,m})_i\models\{\,S'\mid\form'\,\}},
\]
which is Definition~\ref{def:il-valid}.
\end{proof}

\end{document}
